\section{Transport of the primary Frobenius structure}
\label{sec:primary-frobenius}

Section~\ref{sec:gkt-decoration} compares wall-crossing groups.  A comparison
of Frobenius structures requires further data: a primitive form, oscillatory
periods, a residue pairing, and a Kodaira--Spencer basis.  This section
supplies these data on the primary slice, where every descendent variable
vanishes, starting from the primary potential of \cite{grossetal2022open}
with the primitive form $dx\wedge dy$ in the sense of
\cite[Def.~4.3]{grossetal2022open}, recalled below.  The annular geometry
enters only through the transition maps between charts; these are algebra
isomorphisms onto their images, and they preserve the Frobenius structure by
transport of structure (Proposition~\ref{thm:primary-Frobenius}).  The source
structure is calibrated in Appendix~\ref{app:primary-calibration},
independently of the annulus: the primary potential has a unique
representative $W^{\GKT}_{\mathrm{prim}}$ invariant under the exchange of the
coordinates, for which $dx\wedge dy$ is a primitive form with the nine
primary variables as flat coordinates (Lemma~\ref{lem:reflection-fixed} and
the discussion after it).
On the marginal line $W^{\GKT}_{\mathrm{prim}}$ is the flat structure of
Noumi and Yamada with primitive form normalized to $dx\wedge dy$, and its
coefficients give the first open invariants of the reflection-fixed chamber
index (Theorem~\ref{thm:marginal-calibration}).

In the \emph{source chart}, given by the fixed frame of the sheet $S_0$
(Sections~\ref{sec:affine-kummer} and~\ref{sec:jacobi-theta}), we write
$x=z_1$ and $y=z_2$.  Under the dictionary of Section~\ref{sec:jacobi-theta}
these are the coordinates of the GKT potentials in \cite{grossetal2022open}
and in Sections~\ref{sec:gkt-decoration} and~\ref{sec:positive-descendents},
not the Cox coordinates of \eqref{eq:quartic-pencil}.  Then
$dx\wedge dy=\Omega_\kappa$ by \eqref{eq:primitive-form}, and the
Hamiltonian fields $X_{a,b}=X^\kappa_{\kappa+(a,b)}$ of
Proposition~\ref{prop:shifted-bracket}, of divergence zero for
$dx\wedge dy$, are
\begin{equation}\label{eq:quartic-vector-field}
 X_{a,b}
 =x^ay^b\bigl((b+1)x\partial_x-(a+1)y\partial_y\bigr).
\end{equation}
The symmetric deformation variables of \cite[Def.~4.10]{grossetal2022open}
are $t_{a,b,d}$, $0\leq a,b\leq2$, $d\geq0$, where $(a,b)$ is a twist and $d$
is the descendent index.  The \emph{primary variables} are those with $d=0$,
and we write
$t_{a,b}=t_{a,b,0}$.  The variable $t_{a,b}$ is the parameter dual to the
monomial $x^ay^b$; the nine monomials $x^ay^b$, $0\leq a,b\leq2$, form a basis
of the Jacobian ring $\Jac(x^4+y^4)$ and correspond to the nine narrow
elements $(e^{2\pi i(a+1)/4},e^{2\pi i(b+1)/4})$ of $\mu_4\times\mu_4$,
those whose fixed locus in $\CC^2$ is the origin.  Put
\begin{equation}\label{eq:primary-base}
 A_{\mathrm{prim}}
 =\QQ[[t_{a,b}:0\leq a,b\leq2]],
 \qquad \mathfrak m_{\mathrm{prim}}=(t_{a,b}).
\end{equation}
A statement over this complete ring means a compatible system of statements
over the quotients $\QQ[t_{a,b}]/(t_{a,b}^{n_{a,b}+1})$.  For the source
potential the compatibility is part of Lemma~\ref{lem:reflection-fixed}.
The exchange of the coordinates and the degree-zero field, the case
$(a,b)=(0,0)$ of \eqref{eq:quartic-vector-field}, are
\begin{equation}\label{eq:reflection}
 \sigma(x,y,t_{a,b})=(y,x,t_{b,a}),
 \qquad
 X_{0,0}=x\partial_x-y\partial_y .
\end{equation}

The periods are taken over the good basis of
\cite[\S4.1]{grossetal2022open}: the product cycles $\Xi_{a,b}$,
$0\leq a,b\leq2$, in the relative homology
$H_2(\CC^2,\operatorname{Re}(x^4+y^4)/\hbar\ll0;\CC)$, dual to the
monomial basis in the sense that
\[
 \int_{\Xi_{a,b}}x^cy^d\,e^{(x^4+y^4)/\hbar}\,dx\wedge dy
 =\delta_{ac}\delta_{bd},
 \qquad 0\leq c,d\leq2
\]
\cite[(4.3)]{grossetal2022open}.  The same integrals for arbitrary
$c,d\geq0$, the \emph{moments}, follow by integration by parts
\cite[Lem.~4.2]{grossetal2022open}.  For a deformation $W$ of $x^4+y^4$ over
a complete local ring, the period of $(W,dx\wedge dy)$ over $\Xi_{a,b}$ is
the formal series in $\hbar$ and $\hbar^{-1}$ obtained by expanding
$\int_{\Xi_{a,b}}e^{W/\hbar}\,dx\wedge dy$ in the moments.  Following
\cite[Def.~4.3]{grossetal2022open}, $dx\wedge dy$ is a \emph{primitive form}
for $W$ if, for every $0\leq a,b\leq2$, the period over $\Xi_{a,b}$ contains
no positive power of $\hbar$ and has constant term
$\delta_{(a,b),(0,0)}$; the coefficients of $\hbar^{-1}$ in the nine periods
are then the \emph{flat coordinates}.  Thus $dx\wedge dy$ is a primitive
form with flat coordinates $t_{a,b}$ exactly when
\[
 \int_{\Xi_{a,b}}e^{W/\hbar}\,dx\wedge dy
 =\delta_{(a,b),(0,0)}+\hbar^{-1}t_{a,b}+O(\hbar^{-2}),
 \qquad 0\leq a,b\leq2,
\]
which is the form in which \cite[Cor.~4.17]{grossetal2022open} states it.
GKT formulate the definition for a form $f\,dx\wedge dy$ and present it as
a simplified version of Saito's theory of primitive forms; only $f=1$ occurs
here.

\subsection{The source Frobenius algebra}
\label{subsec:source-Frobenius}

At the Fermat point the primary state space is
\begin{equation}\label{eq:Fermat-Jacobian}
 H_{4,4}=\QQ[x,y]/(x^3,y^3),
 \qquad e_{a,b}=[x^ay^b],\quad 0\leq a,b\leq2.
\end{equation}
Normalize the Grothendieck residue trace by
\begin{equation}\label{eq:residue-trace}
 \epsilon_0(f)
 =16\operatorname{Res}_0
   \frac{f\,dx\wedge dy}{(4x^3)(4y^3)}.
\end{equation}

\begin{proposition}[Residue pairing at the Fermat point]
\label{prop:classical-vertex}
The pairing $\eta_0(f,g)=\epsilon_0(fg)$ on $H_{4,4}$ is nondegenerate and
\begin{equation}\label{eq:classical-pairing}
 \eta_0(e_{a,b},e_{c,d})
 =\delta_{a+c,2}\delta_{b+d,2}.
\end{equation}
Equivalently, $\epsilon_0$ is the coefficient of the socle monomial
$e_{2,2}$, and $e_{a,b}e_{2-a,2-b}=e_{2,2}$ for $0\leq a,b\leq2$.
\end{proposition}

\begin{proof}
The residue in \eqref{eq:residue-trace} extracts the coefficient of
$x^2y^2$.  Since $e_{a,b}e_{c,d}=e_{a+c,b+d}$ when $a+c\leq2$ and $b+d\leq2$
and vanishes otherwise, this proves \eqref{eq:classical-pairing}, the socle
identity, and nondegeneracy.
\end{proof}

Over the primary base, the Jacobian algebra of the reflection-fixed primary
potential and its normalized residue trace are
\begin{equation}\label{eq:source-primary-Jacobian}
 \mathcal R_0=A_{\mathrm{prim}}[x,y]\big/
  (\partial_xW^{\GKT}_{\mathrm{prim}},
   \partial_yW^{\GKT}_{\mathrm{prim}}),
 \qquad
 \epsilon(h)=16\operatorname{Res}_0
 \frac{h\,dx\wedge dy}
 {\partial_xW^{\GKT}_{\mathrm{prim}}\,\partial_yW^{\GKT}_{\mathrm{prim}}},
\end{equation}
with residue pairing $\eta(f,g)=\epsilon(fg)$ and Kodaira--Spencer classes
$\Phi^0_{a,b}=[\partial W^{\GKT}_{\mathrm{prim}}/\partial t_{a,b}]$,
$0\leq a,b\leq2$.

\begin{lemma}[Source Frobenius algebra]
\label{lem:source-Frobenius}
The classes $\Phi^0_{a,b}$ form an $A_{\mathrm{prim}}$-basis of
$\mathcal R_0$, with $\Phi^0_{a,b}=x^ay^b+O(\mathfrak m_{\mathrm{prim}})$
and $\Phi^0_{0,0}=[1]$.  The pairing $\eta$ is nondegenerate and makes
$\mathcal R_0$ a Frobenius algebra.  Under the Kodaira--Spencer map
$\partial/\partial t_{a,b}\mapsto\Phi^0_{a,b}$, the multiplication of
$\mathcal R_0$ and $\eta$ are the product and the flat metric of the primary
Saito--Givental Frobenius manifold of $W^{\GKT}_{\mathrm{prim}}$; in
particular $\eta$ is constant in the flat coordinates $t_{a,b}$.  By
\cite[Thm.~0.3 and Rem.~0.4(2)]{grossetal2022open}, this Frobenius manifold
is the genus-zero primary closed FJRW Frobenius manifold of
$(x^4+y^4,\mu_4^2)$.
\end{lemma}

\begin{proof}
Modulo quadratic terms in the $t_{a,b}$,
$W^{\GKT}_{\mathrm{prim}}\equiv x^4+y^4+\sum t_{a,b}x^ay^b$
\cite[Def.~3.47(1) and (4.11)]{grossetal2022open}, which gives the
congruence for $\Phi^0_{a,b}$.  Over each quotient $A^{\mathrm{prim}}_I$,
where $\mathfrak m_{\mathrm{prim}}$ is nilpotent, the nine monomials $x^ay^b$
span $\mathcal R_0$ by Nakayama's lemma, since they span
$\mathcal R_0/\mathfrak m_{\mathrm{prim}}\mathcal R_0=H_{4,4}$.  The
functional $\epsilon$ vanishes on the Jacobian ideal, and the Gram matrix of
$\eta$ on these monomials reduces modulo $\mathfrak m_{\mathrm{prim}}$ to the
invertible matrix of Proposition~\ref{prop:classical-vertex}, so its
determinant is a unit.  Hence $\mathcal R_0$ is free with basis the
$x^ay^b$ and $\eta$ is nondegenerate; the $\Phi^0_{a,b}$ form a basis by
Nakayama's lemma, $\Phi^0_{0,0}=[1]$ by \eqref{eq:constant-deformation}, and
$\eta(fg,h)=\epsilon(fgh)=\eta(f,gh)$.  The flat coordinates $t_{a,b}$ of
$(W^{\GKT}_{\mathrm{prim}},dx\wedge dy)$ (the discussion after
Lemma~\ref{lem:reflection-fixed}) agree with those of Saito--Givental theory
in the good-basis description of \cite[\S2.2]{heetal2022}
\cite[Rems.~0.4(2) and~4.5]{grossetal2022open}, so the Kodaira--Spencer map
identifies the product and flat metric of Saito's Frobenius manifold
\cite{saito1983period} with the multiplication of $\mathcal R_0$ and $\eta$.
Every symmetric chamber index is enumerative
\cite[Prop.~3.46 and Cor.~5.14]{grossetal2022open}, so
\cite[Thm.~0.3 and Rem.~0.4(2)]{grossetal2022open} applies.
\end{proof}

\subsection{Transport through the annular charts}
\label{subsec:KS-classes}

Let $R$ be an Artinian local $\QQ$-algebra of labelled deformation
coefficients, base-changed to $\Bbbk$ where it meets the chart algebras
(Section~\ref{sec:introduction}).  For $W\in R[x,y]$ the Brieskorn module is
\begin{equation}\label{eq:Brieskorn-module}
 \operatorname{GM}_\hbar(W)
 =H^2\!\left(
  \Omega_{R[x,y]/R}^{\bullet}[[\hbar]],
  \hbar d+dW\wedge
 \right).
\end{equation}
For an admissible path $\mathfrak p$ (Definition~\ref{def:framed-path}), let
$P_{\mathfrak p}$ be the composite of the transition maps along it: the slab
maps \eqref{eq:local-root-ring} with the affine holonomy they carry, and the
wall automorphisms.  These are injective homomorphisms of $R$-algebras
between chart algebras, so $U_{\mathfrak p}=P_{\mathfrak p}(x)$ and
$V_{\mathfrak p}=P_{\mathfrak p}(y)$ are algebraically independent over $R$,
and $P_{\mathfrak p}$ is an isomorphism onto the image algebra
\begin{equation}\label{eq:image-algebra}
 P_{\mathfrak p}\bigl(R[x,y]\bigr)=R[U_{\mathfrak p},V_{\mathfrak p}]
\end{equation}
inside the target chart algebra.  Brieskorn modules, Jacobian algebras,
residue pairings and good-basis periods (the moment functionals of
\cite[Lem.~4.2]{grossetal2022open}) of $P_{\mathfrak p}W$ are formed over
this image algebra by the same formulas in $U_{\mathfrak p},V_{\mathfrak p}$,
not over the target chart algebra, a Laurent algebra over which the Jacobian
quotient of a transported primary potential vanishes because its
$U_{\mathfrak p}$-derivative is a unit.  The periods are
formal functionals, since the torus of the target chart does not contain the
critical point $U_{\mathfrak p}=V_{\mathfrak p}=0$.  By \eqref{eq:slab-conformal}, $dU_{\mathfrak p}\wedge
dV_{\mathfrak p}$ is the form $\Omega'_\kappa$ of the target chart times the
unit recorded by the line $\Lcal_\kappa$.

On the primary slice put
$W_{\mathfrak p}=P_{\mathfrak p}(W^{\GKT}_{\mathrm{prim}})$,
\begin{equation}\label{eq:transported-primary-Jacobian}
 \mathcal R_{\mathfrak p}
 =A_{\mathrm{prim}}[U_{\mathfrak p},V_{\mathfrak p}]\big/
  (\partial_{U_{\mathfrak p}}W_{\mathfrak p},
   \partial_{V_{\mathfrak p}}W_{\mathfrak p}),
\end{equation}
let $\eta_{\mathfrak p}$ be the residue pairing of
$(W_{\mathfrak p},dU_{\mathfrak p}\wedge dV_{\mathfrak p})$, and let
\begin{equation}\label{eq:KS-sections}
 \Phi_{a,b}^{\mathfrak p}
 =\left[\frac{\partial W_{\mathfrak p}}{\partial t_{a,b}}\right]
 \in\mathcal R_{\mathfrak p},
 \qquad 0\leq a,b\leq2,
\end{equation}
be the transported Kodaira--Spencer classes; for the trivial path these are
the $\Phi^0_{a,b}$ of Lemma~\ref{lem:source-Frobenius}.

\begin{proposition}[Transported primary Frobenius structure]
\label{thm:primary-Frobenius}
\label{prop:Brieskorn-transport}
\label{thm:primary-Saito-system}
Let $\mathfrak p$ be an admissible path.
\begin{enumerate}
\item For $W\in R[x,y]$, $P_{\mathfrak p}$ induces an isomorphism of twisted
de Rham complexes over $R[[\hbar]]$, and hence
\begin{equation}\label{eq:Brieskorn-transport}
 \operatorname{GM}_\hbar(W)
 \xrightarrow{\sim}
 \operatorname{GM}_\hbar(P_{\mathfrak p} W),
\end{equation}
carrying $[dx\wedge dy]$ to $[dU_{\mathfrak p}\wedge dV_{\mathfrak p}]$.
Modulo $\hbar$ it identifies the Jacobian algebras and the residue pairings of
$(W,dx\wedge dy)$ and $(P_{\mathfrak p}W,dU_{\mathfrak p}\wedge
dV_{\mathfrak p})$, and it preserves every good-basis period.
\item On the primary slice, $P_{\mathfrak p}$ is the base change to
$A_{\mathrm{prim}}$ of an injective homomorphism of chart algebras over
$\Bbbk$, independent of the $t_{a,b}$.  For the transported periods,
$dU_{\mathfrak p}\wedge dV_{\mathfrak p}$ is a primitive form for
$W_{\mathfrak p}$ in the sense of \cite[Def.~4.3]{grossetal2022open}, and the
nine $t_{a,b}$ are simultaneous flat coordinates.
\item The classes \eqref{eq:KS-sections} form an $A_{\mathrm{prim}}$-basis of
$\mathcal R_{\mathfrak p}$ with
\begin{equation}\label{eq:KS-naturality}
 \Phi_{a,b}^{\mathfrak p}=P_{\mathfrak p}(\Phi_{a,b}^0),
 \qquad
 \Phi_{a,b}^0=x^ay^b+O(\mathfrak m_{\mathrm{prim}}),
\end{equation}
and their multiplication, the unit $\Phi_{0,0}^{\mathfrak p}=[1]$ and the
residue pairing $\eta_{\mathfrak p}$ have coefficient matrices independent of
$\mathfrak p$, namely the structure constants of the primary Saito--Givental
Frobenius manifold of $W^{\GKT}_{\mathrm{prim}}$, which
\cite[Thm.~0.3 and Rem.~0.4(2)]{grossetal2022open} identifies with the
genus-zero primary closed FJRW Frobenius manifold of $(x^4+y^4,\mu_4^2)$.
\end{enumerate}
These maps compose under concatenation of paths and are compatible with
reduction of $R$, in particular to the quotients $A^{\mathrm{prim}}_I$, and
with refinement of the root-stack presentation.  For a second admissible path
$\mathfrak p'$, the isomorphism
$P_{\mathfrak p'}P_{\mathfrak p}^{-1}:R[U_{\mathfrak p},V_{\mathfrak p}]\to
R[U_{\mathfrak p'},V_{\mathfrak p'}]$ has the properties in (1) and carries
$\Phi^{\mathfrak p}_{a,b}$ to $\Phi^{\mathfrak p'}_{a,b}$.
\end{proposition}

\begin{proof}
Each part is transport of structure along $P_{\mathfrak p}$.  (1) Since
$P_{\mathfrak p}$ carries $dx,dy$ to $dU_{\mathfrak p},dV_{\mathfrak p}$, it
satisfies $(\hbar d+d(P_{\mathfrak p}W)\wedge)P_{\mathfrak p}
=P_{\mathfrak p}(\hbar d+dW\wedge)$ on forms, so pullback is a chain
isomorphism; modulo $\hbar$ it is an isomorphism of Koszul complexes, whose
top cohomology is the Jacobian algebra times the volume form.  The residue
and the moment functionals are given by the same formulas in both coordinate
systems, so they are preserved.  Composition, reduction and refinement
follow from functoriality of the transition maps, and the argument applies
equally to $P_{\mathfrak p'}P_{\mathfrak p}^{-1}$.

(2) Every wall automorphism is the identity modulo $\mathfrak m=(t_1,t_2)$
(Definition~\ref{def:Jacobi-wall}), an ideal generated by descendent
variables, so on the primary slice $P_{\mathfrak p}$ is a composite of slab
and holonomy maps, which involve no deformation parameters.  The conditions
of \cite[Def.~4.3]{grossetal2022open} concern only the good-basis periods; they hold for $(W^{\GKT}_{\mathrm{prim}},dx\wedge dy)$ by
the discussion after Lemma~\ref{lem:reflection-fixed} and are preserved by
(1).

(3) By the chain rule and (2), $P_{\mathfrak p}$ induces an isomorphism of
$A_{\mathrm{prim}}$-algebras $\mathcal R_0\to\mathcal R_{\mathfrak p}$ that
commutes with $\partial/\partial t_{a,b}$, which gives the first identity in
\eqref{eq:KS-naturality}, carries $[1]$ to $[1]$, and by (1) carries $\eta$
to $\eta_{\mathfrak p}$.  Hence the multiplication, the unit and the pairing
have the same coefficient matrices in the bases $\Phi^{\mathfrak p}_{a,b}$ as
for the trivial path, $P_{\mathfrak p'}P_{\mathfrak p}^{-1}$ carries
$\Phi^{\mathfrak p}_{a,b}$ to $\Phi^{\mathfrak p'}_{a,b}$, and the remaining
assertions follow from Lemma~\ref{lem:source-Frobenius}.
\end{proof}

Wall factors, and the residual factors by which they differ from the GKT
factors (Section~\ref{sec:gkt-decoration}), are flows of the fields
\eqref{eq:quartic-vector-field} with nilpotent coefficients, hence
volume-preserving right equivalences of $(W,dx\wedge dy)$.  By
Proposition~\ref{thm:primary-Frobenius}(1), the identities between algebra
automorphisms of Appendix~\ref{app:residual-coherence} therefore become
composition identities of Brieskorn chain maps.  Since $P_{\mathfrak p}$ carries Jacobian ideals to
Jacobian ideals, a coefficient of the Jacobian algebra on transported source
sections, such as $1/32$ in \eqref{eq:marginal-product}, does not depend on
the admissible path, although $U_{\mathfrak p}$ and $V_{\mathfrak p}$ do,
through the hyperbolic holonomy.

\begin{remark}[Deck transformations and core holonomy]
\label{rem:reflection-deck}
The deck transformations of the fractional-exponent cover $p_{\mathrm{fr}}$
and of the twist cover, and the core holonomy, relate admissible paths
$\mathfrak p,\mathfrak p'$ ending over the same region of $\mathfrak D$.  By
Proposition~\ref{thm:primary-Frobenius},
$\phi=P_{\mathfrak p'}P_{\mathfrak p}^{-1}:\mathcal R_{\mathfrak p}\to\mathcal
R_{\mathfrak p'}$ is $A_{\mathrm{prim}}$-linear and carries transported bases
to transported bases, so the structure constants are invariant under these
maps.  If the two sheets of $p_{\mathrm{fr}}$ over a region are identified by
such a $\phi$, the involution $(r,r')\mapsto(\phi^{-1}r',\phi r)$ splits
$\mathcal R_{\mathfrak p}\oplus\mathcal R_{\mathfrak p'}$ into its
$(\pm1)$-eigensummands, free of rank nine over $A_{\mathrm{prim}}$ and
orthogonal for the residue pairing; the $(+1)$-eigensummand
$\{(r,\phi r)\}$ is a unital subalgebra isomorphic to $\mathcal R_{\mathfrak
p}$.
\end{remark}
